% Part: computability
% Chapter: machines-computations
% Section: halting-problem

\documentclass[../../../include/open-logic-section]{subfiles}

\begin{document}

\olfileid{tur}{und}{hal}
\olsection{ఆగే సమస్య}

\begin{explain}
ట్యూరింగ్ యంత్ర వర్ణనల ఏదో ఒక లెక్కింపును స్థిరపరిచామని అనుకుందాం.
అప్పుడు ప్రతి ట్యూరింగ్ యంత్రానికి ఒక \emph{సూచిక} వస్తుంది: ట్యూరింగ్ యంత్ర
వర్ణనల $M_1$, $M_2$, $M_3$, \dots{} లెక్కింపులో దాని స్థానం.

ట్యూరింగ్-గణనీయాలు కాని ప్రమేయాలు తప్పక ఉన్నాయని మనకు తెలుసు: ట్యూరింగ్
యంత్ర వర్ణనల సమితి---అందువల్ల ట్యూరింగ్ యంత్రాల సమితి---\tetoken{లెక్కించదగిన}{!!{enumerable}}; కానీ
$\Nat$ నుండి $\Nat$కు వెళ్లే అన్ని ప్రమేయాల సమితి అలా కాదు. అయినా గణనీయాలు
కాని ప్రమేయాలకు నిర్దిష్ట ఉదాహరణలను కూడా కనుగొనవచ్చు. అలాంటి ప్రమేయం ఒకటి
ఆగే ప్రమేయం.
\end{explain}

\begin{defn}[ఆగే ప్రమేయం]
 \emph{ఆగే ప్రమేయం}~$h$ను ఇలా నిర్వచిస్తాం:
\[
h(e,n) =
\begin{cases}
  \text{0} & \text{ఇన్‌పుట్ $n$కు యంత్రం~$M_e$ ఆగకపోతే} \\
  \text{1} & \text{ఇన్‌పుట్ $n$కు యంత్రం~$M_e$ ఆగితే}
\end{cases}
\]
\end{defn}

\begin{defn}[ఆగే సమస్య]
ఏ $e$, $n$కైనా, $n$ గీతల ఇన్‌పుట్‌పై ట్యూరింగ్ యంత్రం~$M_e$ ఆగుతుందో లేదో
నిర్ణయించే సమస్యను \emph{ఆగే సమస్య} అంటారు.
\end{defn}

\begin{explain}
సంబంధిత ప్రమేయం~$s$ ట్యూరింగ్-గణనీయం కాదని చూపడం ద్వారా~$h$
ట్యూరింగ్-గణనీయం కాదని చూపుతాం. ట్యూరింగ్ యంత్రం గణించగల దేన్నైనా ఒక
క్రమశిక్షిత ట్యూరింగ్ యంత్రం గణించగలదనే వాస్తవం (\olref[mac][dis]{sec}),
రెండు ట్యూరింగ్ యంత్రాలను కలిపి ఒకే యంత్రాన్ని నిర్మించవచ్చనే వాస్తవం
(\olref[mac][cmb]{sec}) ఈ నిరూపణకు ఆధారాలు.
\end{explain}

\begin{defn} ప్రమేయం~$s$ను ఇలా నిర్వచిస్తాం:
\[
s(e) =
\begin{cases}
  \text{0} & \text{ఇన్‌పుట్ $e$కు యంత్రం~$M_e$ ఆగకపోతే} \\
  \text{1} & \text{ఇన్‌పుట్ $e$కు యంత్రం~$M_e$ ఆగితే}
\end{cases}
\]
\end{defn}

\begin{lem}
ప్రమేయం~$s$ ట్యూరింగ్-గణనీయం కాదు.
\end{lem}

\begin{proof}
వైరుధ్యం కోసం ప్రమేయం~$s$ ట్యూరింగ్-గణనీయం అని అనుకుందాం. అప్పుడు~$s$ను
గణించే ట్యూరింగ్ యంత్రం~$S$ ఉండేది. $S$ ఆగినప్పుడు మొదటి గడిని చదువుతూనే
ఆగుతుందని (అంటే, అది క్రమశిక్షితమని) సామాన్యతకు భంగం లేకుండా అనుకోవచ్చు.
ఈ యంత్రాన్ని మరొక యంత్రం~$J$తో ``కలపవచ్చు.'' దాన్ని ఇన్‌పుట్~$0$పై
ప్రారంభిస్తే అది ఆగుతుంది (అంటే, ప్రారంభ స్థితిలో టేపు-అంత సంకేతానికి
కుడివైపు గడిని చదువుతున్నప్పుడు $\TMblank$ను చదివితే ఆగుతుంది); లేకపోతే
అది కుడివైపుకు వెళ్లిపోతూ ఎప్పటికీ ఆగదు. $S$ను $J$తో కలిపి నిర్మించిన
యంత్రం $S \concat J$ ఒక ట్యూరింగ్ యంత్రం; కాబట్టి ఏదో ఒక~$e$కు అది $M_e$
(అంటే, లెక్కింపులో ఎక్కడో ఒకచోట కనిపిస్తుంది). $M_e$ను $e$~$\TMstroke$ల
ఇన్‌పుట్‌పై ప్రారంభించండి. రెండు అవకాశాలు ఉన్నాయి: $M_e$ ఆగుతుంది లేదా ఆగదు.
\begin{enumerate}
\item $e$~$\TMstroke$ల ఇన్‌పుట్‌కు $M_e$ ఆగుతుందని అనుకుందాం. అప్పుడు
  $s(e) = 1$. కాబట్టి $S$ను~$e$పై ప్రారంభించినప్పుడు, టేపుపై ఒకే
  $\TMstroke$ను అవుట్‌పుట్‌గా ఇచ్చి ఆగుతుంది. అప్పుడు $J$ టేపుపై ఒక
  $\TMstroke$తో మొదలవుతుంది. ఆ సందర్భంలో $J$ ఆగదు. కానీ $M_e$ అనేది
  $S \concat J$ యంత్రం; కాబట్టి $S$, తరువాత $J$ చేసే పనినే అది కచ్చితంగా
  చేయాలి (అంటే ఈ సందర్భంలో కుడివైపుకు వెళ్లిపోతూ ఎప్పటికీ ఆగకూడదు).
  అందువల్ల $e$~$\TMstroke$ల ఇన్‌పుట్‌కు $M_e$ ఆగలేదు.

\item ఇప్పుడు $e$~$\TMstroke$ల ఇన్‌పుట్‌కు $M_e$ ఆగదని అనుకుందాం.
  అప్పుడు $s(e) = 0$; $S$ను ఇన్‌పుట్~$e$పై ప్రారంభించినప్పుడు ఖాళీ టేపుతో
  ఆగుతుంది. ఖాళీ టేపుపై ప్రారంభించిన~$J$ వెంటనే ఆగుతుంది. మళ్లీ, $S$,
  తరువాత~$J$ చేసే పనినే $M_e$ చేస్తుంది; కాబట్టి $e$~$\TMstroke$ల
  ఇన్‌పుట్‌కు $M_e$ తప్పక ఆగాలి.
\end{enumerate}
ప్రతి సందర్భంలోనూ మన పరికల్పనకు వైరుధ్యం వస్తుంది. ట్యూరింగ్ యంత్రం~$S$
ఉండదని ఇది చూపుతుంది: $s$ ట్యూరింగ్-గణనీయం కాదు.
\end{proof}

\begin{thm}[ఆగే సమస్య యొక్క అపరిష్కార్యత]
\ollabel{thm:halting-problem} ఆగే సమస్య పరిష్కరించలేనిది; అంటే, ప్రమేయం~$h$
ట్యూరింగ్-గణనీయం కాదు.
\end{thm}

\begin{proof}
$h$ను ట్యూరింగ్ యంత్రం~$H$ గణిస్తుందని అనుకుందాం. అప్పుడు~$s$ను గణించే
ట్యూరింగ్ యంత్రాన్ని $H$తో నిర్మించవచ్చు: మొదట ఇన్‌పుట్‌కు ఒక నకలు చేసి
(ఒక~$\TMblank$ సంకేతంతో వేరుచేసి), తరువాత మొదటికి తిరిగి వెళ్లి~$H$ను
నడపాలి. మొదటి పని చేసే యంత్రాన్ని స్పష్టంగా నిర్మించవచ్చు
(\cref{tur:mac:dis:prob:copier}); $H$ ఉంటే, దాన్ని అలాంటి నకలు యంత్రంతో
``కలిపి,'' $M_e$ ఇన్‌పుట్~$e$పై ఆగుతుందో లేదో నిర్ణయించే, అంటే~$s$ను గణించే,
కొత్త యంత్రాన్ని పొందగలిగేవారం. అలాంటి యంత్రం ఉండదని ఇప్పటికే చూపాం.
అందువల్ల~$h$ కూడా ట్యూరింగ్-గణనీయం కాదు.
\end{proof}

\begin{prob}
మూడు-ఆగే సమస్య అంటే, ఏకపక్షంగా ఎంచుకున్న ట్యూరింగ్ యంత్రం, మిగతా
టేపంతా ఖాళీగా ఉండి మూడు $\TMstroke$లు ఇన్‌పుట్‌గా ఉన్నప్పుడు ఆగుతుందో లేదో
నిర్ణయించే నిర్ణయ ప్రక్రియను ఇచ్చే సమస్య. మూడు-ఆగే సమస్య పరిష్కరించలేనిదని
నిరూపించండి.
\end{prob}

\begin{prob}
ట్యూరింగ్ యంత్రం--ఇన్‌పుట్ యుగ్మాలు $M_e$,~$n$లో $e \neq n$ అయినవాటికి ఆగే
సమస్య పరిష్కరించదగినదైతే, $e = n$ అయిన సందర్భాలకు కూడా అది పరిష్కరించదగినదని
చూపండి.
\end{prob}

\begin{prob}
ఇన్‌పుట్ ఒక సంఖ్య~$e$ అయినప్పుడు ఆగే సమస్య పరిష్కరించలేనిదని నిరూపించాం;
అన్ని ట్యూరింగ్ యంత్రాల లెక్కింపు ద్వారా ఈ సంఖ్య ట్యూరింగ్ యంత్రం~$M_e$ను
గుర్తిస్తుంది. దానికి బదులుగా \olref[tur][und][enu]{sec}లోని ట్యూరింగ్ యంత్ర
వర్ణనలనే నేరుగా ఇన్‌పుట్‌గా అనుమతిస్తే ఏమవుతుంది? సూచికలకు బదులుగా ట్యూరింగ్
యంత్ర వర్ణనలను ఇన్‌పుట్‌గా తీసుకుని ఆగే సమస్యను నిర్ణయించే ట్యూరింగ్ యంత్రం
ఉండగలదా? ఎందుకు ఉండగలదో లేదా ఉండలేదో వివరించండి.
\end{prob}

\begin{prob} కింది విధంగా నిర్వచించిన \emph{పాక్షిక} ప్రమేయం~$s'$ ట్యూరింగ్-గణనీయం
అని చూపండి:
  \[
  s'(e) =
  \begin{cases}
    \text{1} & \text{ఇన్‌పుట్ $e$కు యంత్రం~$M_e$ ఆగితే}\\
    \text{నిర్వచితం కాదు} & \text{ఇన్‌పుట్ $e$కు యంత్రం~$M_e$ ఆగకపోతే}
  \end{cases}
  \]
\end{prob}

\end{document}
